\section{De la U-Net gigante a la DiT escalable}

El auge del Transformer no se debió a que la U-Net no pudiera generar buenas imágenes, sino a que, cuando la U-Net evolucionó hacia una gran cantidad de estructuras personalizadas de convolución, residuales, atención, submuestreo y sobremuestreo y skip, se volvió difícil que heredara directamente el escalamiento y las optimizaciones de sistema de los Transformer de uso general. El valor de DiT reside, en primer lugar, en la unificación de la estructura y la reutilización del ecosistema.

\subsection{Por qué la U-Net se volvió cada vez más «gigante»}

Los primeros DDPM y la latent diffusion usaban principalmente U-Net. Su sesgo inductivo multiescala es adecuado para imágenes, pero la U-Net moderna de text-to-image ya no es una simple red convolucional simétrica, sino una mezcla de muchos componentes especializados. Esta sección sigue la pregunta de «por qué un prior de imagen maduro se convierte en una carga de ingeniería»: primero se examinan las capacidades que aporta la ruta multiescala y después se rastrea cómo los block especializados, las interfaces de condicionamiento y la configuración por resolución acumulan complejidad capa por capa.

\lecturefigure{slide-15.jpg}{En la época de DDPM/LDM dominaba la U-Net, que después incorporó gradualmente Transformer blocks}{15}

La U-Net obtiene un campo receptivo amplio mediante el submuestreo, recupera la resolución mediante el sobremuestreo y conserva los detalles con las skip connection. El condicionamiento de texto suele inyectarse en las capas intermedias mediante cross-attention. Su desempeño es maduro, pero hace que cada etapa de resolución tenga distintos canales, número de block y configuración de atención.

\lecturefigure{slide-16.jpg}{Componentes de la U-Net gigante: stem convolucional, downblocks, midblock, upblocks y cabeza de salida}{16}

La convolución de entrada proyecta el noisy latent a los hidden channels; los downblocks reducen gradualmente el tamaño espacial; el midblock procesa la representación más comprimida; los upblocks combinan las skip features para recuperar los detalles. Cada block puede incluir además ResNet, normalization, self/cross-attention y operadores de muestreo.

\lecturefigure{slide-17.jpg}{Los problemas de ingeniería de la «Giant U-Net»: una configuración enorme y difícil de mantener unificada}{17}

El objetivo didáctico de la captura de la configuración no es burlarse de la cantidad de parámetros, sino mostrar la architecture surface area: los arreglos de canales, los tipos de block, las attention head, la cross-attention dimension, el número de capas y el método de muestreo deben coordinarse entre sí. Es fácil que un experimento modifique varias variables a la vez, y la reproducción y la optimización también se vuelven más difíciles.

\lecturefigure{slide-18.jpg}{El interior de un downblock se compone de diversos módulos personalizados de ResNet, Transformer y muestreo}{18}

Los block personalizados pueden aprovechar el prior de imagen, pero también impiden reutilizar directamente optimizaciones estándar como fused kernel, FlashAttention, QK-Norm o GQA. Un fused kernel combina varias operaciones de GPU en un solo kernel para reducir las lecturas y escrituras en HBM (High Bandwidth Memory, memoria de video de alto ancho de banda) y el kernel launch overhead; cuanto más irregular es la estructura, más difícil es la fusión.

\lecturefigure{slide-19.jpg}{La U-Net completa, desplegada, muestra un flujo de datos multiescala complejo y skip connections}{19}

Esta vista panorámica debe leerse siguiendo la resolución: a la izquierda el submuestreo, al centro el bottleneck y a la derecha el sobremuestreo; las flechas horizontales son las skip. La complejidad no proviene solo de la cantidad de parámetros, sino también de la activation memory, de la planificación de tensor de distintos tamaños y de la distribución de las capas de condicionamiento. La optimización del sistema requiere especializarse block por block.

\begin{warningbox}{Complejo no equivale a erróneo, y unificado no equivale a más potente}
El inductive bias multiescala de la U-Net puede seguir siendo valioso con datos limitados o con objetivos de latencia específicos. La estructura de DiT es más ordenada, pero si el número de token es enorme, el costo de la atención supera rápidamente al de la convolución. La elección de arquitectura debe vincularse a la resolución, los datos y el hardware.
\end{warningbox}

\subsection{UViT es una transición; DiT es una reescritura de la estructura}

UViT intenta conservar la jerarquía y las skip largas de la U-Net, al mismo tiempo que sustituye los bloques convolucionales por Transformer. La DiT original va más lejos: aplica patchify al noisy latent y lo envía a un ViT encoder casi estándar, sin mantener ya el tronco multiescala completo en forma de U.

\lecturefigure{slide-20.jpg}{UViT: sustituye los bloques convolucionales por Transformer/MLP y conserva las skip largas entre capas}{20}

UViT muestra que un Transformer puede desempeñar el papel de denoiser, pero su estructura aún conserva rastros de la U-Net. Las skip largas conectan las etapas del codificador y del decodificador, y la longitud de secuencia y la hidden dimension también pueden variar entre capas. Es una «U-Net transformerizada», no «la generación escrita como un modelo estándar de token».

\lecturefigure{slide-21.jpg}{Razones para migrar a DiT: escalamiento, reutilización del ecosistema, sencillez de implementación y calidad de generación}{21}

DiT puede incorporar directamente avances de los Transformer como el MLP en paralelo, QK-Norm, Grouped Query Attention (GQA) y RoPE. GQA hace que varias query heads compartan menos key/value heads para reducir el cómputo y el almacenamiento de KV; QK-Norm, por su parte, normaliza query/key y mejora la estabilidad del entrenamiento de modelos grandes. Lo esencial de la migración es acceder a un ecosistema de backbone de uso general en evolución continua.

\teachervoice{00:14:44--00:21:54, el docente explica «por qué cambiar a DiT» como una unificación de las interfaces de ingeniería e investigación: la U-Net gigante se compone de muchos custom blocks, mientras que DiT puede heredar directamente los avances de la comunidad de Transformer en escalamiento, normalización, atención y kernel.}

\begin{knowledgebox}{Primera aparición: Inductive bias}
El inductive bias (sesgo inductivo) son los supuestos que la arquitectura codifica de antemano. La U-Net enfatiza la convolución local y la multiescala; DiT depende más de que los datos enseñen las relaciones entre token. Cuando aumentan los datos y la capacidad de cómputo, un sesgo manual más débil puede escalar con mayor facilidad, aunque no todas las escalas se benefician.
\end{knowledgebox}

\subsection{Resumen de la sección}

El problema de la U-Net no es que no funcione, sino que las implementaciones modernas están muy personalizadas y es difícil escalarlas de forma unificada. UViT conserva la estructura jerárquica, mientras que la DiT original reescribe el denoising como un token computation cercano a ViT. El siguiente capítulo explicará, punto por punto, qué condicionamientos y mecanismos de inicialización requiere este «cambio aparentemente pequeño».
